% =====================================================================
\section{Stato dell'arte}
\label{sec:sota}
% =====================================================================

\paragraph{Fondamenti di QKD.}
La letteratura sulla QKD è ormai consolidata. La review di Pirandola
et al.~\cite{pirandola2020} copre i protocolli a variabili discrete (BB84 e
le sue varianti con \emph{decoy state}), quelli a variabili continue e le
dimostrazioni di sicurezza. Da qui emergono i due vincoli che guidano tutta
l'ingegneria di rete: il tasso di chiave decresce esponenzialmente con la
distanza e un singolo collegamento in fibra commerciale non supera le poche
centinaia di chilometri.

\paragraph{Dal collegamento alla rete.}
Mehic et al.~\cite{mehic2020} e Cao et al.~\cite{cao2022} descrivono il
passaggio dal collegamento punto-punto alla \emph{rete} QKD: architettura a
strati, \emph{trusted relay}, gestione e instradamento delle chiavi. Il
white paper GÉANT~\cite{geant2024} offre una sintesi leggibile dello stesso
quadro e degli standard ETSI e ITU-T. Il problema dell'instradamento, che
diventa centrale quando si orchestra una rete, è trattato
in~\cite{routing2024}.

\paragraph{Standardizzazione.}
Due organismi hanno prodotto specifiche complementari. L'ETSI ISG QKD ha
iniziato nel 2010 con i casi d'uso~\cite{qkd002} e ha poi definito le
interfacce verso le applicazioni~\cite{qkd004,qkd014}, verso il controller
SDN~\cite{qkd015}, verso l'orchestrator~\cite{qkd018} e tra sistemi di
gestione delle chiavi~\cite{qkd020}. L'ITU-T, con la serie Y.3800, ha
definito l'architettura di riferimento delle reti QKD
(QKDN)~\cite{ituy3800}, il loro controllo e la loro
gestione~\cite{ituy3804} e il controllo tramite SDN~\cite{ituy3805}.

\paragraph{SDN applicato alla QKD.}
Il paradigma \emph{Software-Defined Networking}~\cite{kreutz2015} separa il
piano di controllo dal piano dati e accentra le decisioni in un controller
programmabile. Il gruppo UPM/Telefónica ha mostrato per primo
l'integrazione di risorse QKD in una rete ottica controllata via SDN, sia
sperimentalmente~\cite{aguado2016} sia in un tutorial
sistematico~\cite{aguado2020}; da questo filone nasce la GS~QKD~015.

\paragraph{Orchestrazione e multi-dominio.}
L'interfaccia orchestrator–controller è specificata dalla
GS~QKD~018~\cite{qkd018} e dai relativi moduli YANG~\cite{qkd018yang}. Il
primo caso reale documentato è quello di Sim et al.~\cite{sim2023}
(SK~Telecom), che orchestrano due domini di vendor diversi sul testbed
KOREN. La tesi di Vicario~\cite{vicario2025} realizza un prototipo
gerarchico completo con ONOS e RESTCONF. La rete MadQCI~\cite{madqci2024}
usa estesamente la GS~QKD~015 ma indica l'orchestrazione GS~QKD~018 come
attività ancora in corso. Lavori più recenti propongono applicazioni SDN per
la riconfigurazione dinamica dei link~\cite{mendez2026}, criteri di
selezione dei cammini e instradamento inter-dominio che preserva la
riservatezza della topologia~\cite{laschet2026} e modelli ibridi di
consegna delle chiavi tra reti regionali~\cite{interconnect2026}.

\paragraph{Contributo.}
Rispetto a queste fonti, il presente lavoro offre una lettura unitaria e
ragionata della catena GS~QKD~002 $\to$ 015 $\to$ 018: mostra come i
requisiti dei casi d'uso si traducano nelle componenti di una rete SD-QKD,
spiega la struttura dei modelli YANG che le descrivono e analizza come la
GS~QKD~018 li astragga per l'orchestrazione multi-dominio.
